\chapter{Role-Based Access Control, Networking, and Secrets}
\index{RBAC}\index{access control}\index{roles}\index{network policy}\index{Week 17}%
\begin{objectives}
\begin{itemize}
    \item Apply the system-defined roles (\texttt{ACCOUNTADMIN}, \texttt{SECURITYADMIN}, \texttt{SYSADMIN}, \texttt{USERADMIN}) with separation of duties.
    \item Design custom hierarchies of functional roles and access roles.
    \item Grant privileges only through roles and audit access with \texttt{ACCOUNT\_USAGE} views.
    \item Use secondary roles and session privilege context deliberately.
    \item Build a hands-on production RBAC schema: access roles, functional roles, and user grants.
    \item Architect a least-privilege hierarchy for 500 users across 10 business units, with network policies and secret hygiene.
\end{itemize}
\end{objectives}

\begin{crosscloud}
Identity, role hierarchies, and network perimeters exist in every cloud; Snowflake concentrates all three inside the data platform.

\vspace{2mm}
{\footnotesize
\begin{tabular}{@{}p{2.8cm}p{3.0cm}p{3.0cm}p{3.0cm}@{}}
\toprule
\textbf{Capability} & \textbf{AWS} & \textbf{Azure} & \textbf{GCP} \\
\midrule
Identity system & IAM users/roles & Entra ID & Cloud IAM \\
Privilege grouping & IAM policies \& roles & RBAC role assignments & IAM roles \& groups \\
Separation of duties & Service control policies & Management groups \& PIM & Org policies + IAM \\
Network perimeter & VPC endpoints, SCPs & Private Link, firewall rules & VPC Service Controls \\
Secrets management & Secrets Manager & Key Vault & Secret Manager \\
Audit & CloudTrail & Azure Monitor / Entra logs & Cloud Audit Logs \\
\addlinespace
Snowflake equivalent & \multicolumn{3}{l}{Role hierarchy, network policies, \texttt{ACCOUNT\_USAGE} audit views, key-pair auth} \\
\bottomrule
\end{tabular}}

\vspace{1mm}
\textbf{Microsoft Fabric} layers workspace roles over OneLake security; \textbf{MongoDB} uses database roles and Atlas IP access lists. Snowflake's distinctive rule: \emph{privileges are granted to roles, roles to roles, and roles to users}---never privileges directly to users---which makes access review a graph query instead of an archaeological dig.
\end{crosscloud}

\section{Week 17 Roadmap}
Part IV opens the governance milestone. Every previous chapter created objects; this chapter decides \emph{who may touch them} and proves it to auditors. RBAC is the foundation masking policies (Chapter 18) and secure sharing (Chapter 19) stand on \cite{snowflakeAccessControlDocs}.

Monday covers system roles and administrative separation. Tuesday designs the functional/access role hierarchy. Wednesday addresses secondary roles, network policies, and secrets. Thursday builds a production RBAC schema. Friday scales it to a 500-user enterprise.

\begin{keyterms}
\textbf{Functional role}: a business role describing a job function (ANALYST, ENGINEER), granted to users.\quad
\textbf{Access role}: a technical role owning privileges on specific objects (READ\_GOLD), granted to functional roles.\quad
\textbf{Secondary roles}: the additional role set a session may exercise simultaneously.\quad
\textbf{Network policy}: an IP allow/block list constraining where sessions may originate.\quad
\textbf{Separation of duties}: administrative powers split so no single role controls identity, security, and data.
\end{keyterms}

\section{Monday: System-Defined Roles and Separation of Duties}
\index{ACCOUNTADMIN}\index{SECURITYADMIN}\index{SYSADMIN}\index{USERADMIN}
Snowflake ships a system role tree whose design intent is \emph{separation of duties}:

\begin{table}[H]
\centering
\small
\begin{tabular}{@{}p{3.4cm}p{5.4cm}p{4.6cm}@{}}
\toprule
\textbf{Role} & \textbf{Owns} & \textbf{Must never} \\
\midrule
\texttt{ACCOUNTADMIN} & Account-level settings; inherits everything & Run daily business workloads \\
\texttt{SECURITYADMIN} & Role grants, security objects; manages \texttt{USERADMIN} & Own data objects directly \\
\texttt{SYSADMIN} & Databases, schemas, warehouses & Manage users or grants \\
\texttt{USERADMIN} & Users and custom roles & Touch data objects \\
\texttt{PUBLIC} & Default pseudo-role every user holds & Receive any real privilege \\
\bottomrule
\end{tabular}
\caption{System roles encode separation of duties; respect the boundaries.}
\label{tab:chapter17-system-roles}
\end{table}

The operational rules that follow: daily work happens under custom roles; \texttt{ACCOUNTADMIN} is a break-glass identity with MFA and a tiny membership; and \texttt{PUBLIC} stays empty, because anything granted to it is granted to everyone.

\begin{pitfall}
\textbf{Running pipelines as ACCOUNTADMIN is the most common Snowflake security failure.} It concentrates every permission into the highest-blast-radius identity, defeats audit attribution, and usually persists because ``it worked during setup.'' Week one of any engagement: inventory what is running as \texttt{ACCOUNTADMIN} and demote it.
\end{pitfall}

\section{Tuesday: Functional Roles and Access Roles}
\index{functional role}\index{access role}\index{role hierarchy}
The scalable pattern separates \emph{what a job is} from \emph{what objects it touches}:

\begin{itemize}
    \item \textbf{Access roles} own object privileges: \texttt{READ\_GOLD} has \texttt{USAGE} on the gold database/schema plus \texttt{SELECT} on its tables; \texttt{WRITE\_SILVER} adds \texttt{INSERT}/\texttt{UPDATE} on silver.
    \item \textbf{Functional roles} describe jobs: \texttt{ANALYST}, \texttt{ENGINEER}, \texttt{PIPELINE\_RUNNER}.
    \item \textbf{Grants connect the layers:} access roles to functional roles, functional roles to users. New employee? Grant one functional role. New gold schema? Extend one access role.
\end{itemize}

\begin{lstlisting}[language=SQL, caption={Access roles rolled into functional roles}]
CREATE ROLE READ_GOLD;
GRANT USAGE ON DATABASE analytics TO ROLE READ_GOLD;
GRANT USAGE ON SCHEMA analytics.gold TO ROLE READ_GOLD;
GRANT SELECT ON ALL TABLES IN SCHEMA analytics.gold TO ROLE READ_GOLD;
GRANT SELECT ON FUTURE TABLES IN SCHEMA analytics.gold TO ROLE READ_GOLD;

CREATE ROLE ANALYST;
GRANT ROLE READ_GOLD TO ROLE ANALYST;
GRANT ROLE ANALYST TO USER jane_analyst;
-- object access flows: user -> functional role -> access role -> objects
\end{lstlisting}

Note the \texttt{FUTURE} grant: without it, every new table needs a manual grant---the toil that tempts teams into \texttt{GRANT ... TO PUBLIC}. Future grants make the least-privilege design the \emph{lazy} design, which is the only kind that survives.

\begin{tipbox}
Name roles so the hierarchy reads itself: \texttt{AR\_<domain>\_<level>} for access roles (\texttt{AR\_GOLD\_R}), \texttt{FR\_<function>} for functional roles (\texttt{FR\_DATA\_ANALYST}). Auditors and offboarding scripts both consume role names programmatically.
\end{tipbox}

\section{Wednesday: Secondary Roles, Network Policies, and Secrets}
\index{secondary roles}\index{network policy}\index{secrets}
\subsection{Secondary Roles}
\index{USE SECONDARY ROLES}
A session has a \emph{primary} role (whose privileges create objects and whose name appears in ownership) and may activate \emph{secondary} roles to exercise additional privileges simultaneously---\texttt{USE SECONDARY ROLES ALL} lets a user with \texttt{ANALYST} primary still read through their \texttt{ENGINEER} grants. The sharp edge: \textbf{objects are created under the primary role}. An engineer creating tables with \texttt{ACCOUNTADMIN} primary and secondary roles active produces incorrectly owned objects that the intended roles cannot manage.

\subsection{Network Policies}
\index{network policy}
A network policy constrains \emph{where} sessions may originate: \texttt{ALLOWED\_IP\_LIST} / \texttt{BLOCKED\_IP\_LIST} at account or user level. Combined with MFA and key-pair authentication, network policies close the ``stolen credential used from an unexpected network'' scenario. Service accounts for pipelines should carry the strictest policies---their source IPs (the CI runner fleet, the Airflow host) are known exactly \cite{snowflakeNetworkPoliciesDocs}.

\subsection{Secrets Hygiene}
\index{key-pair authentication}
Snowflake holds no password requirement to use passwords well---but the engineering standard is: human users on SSO with MFA; service users on key-pair authentication with keys in the platform's secrets manager; nothing long-lived in environment files committed anywhere. \texttt{SNOWFLAKE.ACCOUNT\_USAGE.LOGIN\_HISTORY} audits the result: alert on password-based logins for service users and on any login from outside policy networks.

\section{Thursday Hands-On Project: A Production RBAC Schema}
\index{hands-on project!RBAC}
This lab implements the two-layer hierarchy with auditing proof.

\subsection*{Scenario}
The analytics platform needs governed access: analysts read gold, engineers write silver and read gold, pipelines write raw through gold---and security wants a weekly audit query proving it.

\subsection*{Procedure}
\begin{enumerate}
    \item Create access roles \texttt{READ\_GOLD} and \texttt{WRITE\_SILVER} with object and future grants.
    \item Create functional roles \texttt{ANALYST} and \texttt{ENGINEER}; grant the access roles appropriately (\texttt{ENGINEER} also receives \texttt{READ\_GOLD}).
    \item Create users and grant functional roles only---never object privileges directly.
    \item Verify effective access: as an analyst, read gold and fail to write silver; record the expected error.
    \item Write the audit queries: \texttt{GRANTS\_TO\_ROLES} role graph, \texttt{GRANTS\_TO\_USERS}, and the ``who can read table X'' expansion via role inheritance.
    \item Attempt a privilege escalation (analyst self-grant) and capture the denial for the evidence pack.
\end{enumerate}

\subsection*{Deliverables}
\begin{itemize}
    \item \textbf{RBAC DDL script:} roles, grants, future grants, user assignments.
    \item \textbf{Role-hierarchy diagram} generated from \texttt{GRANTS\_TO\_ROLES}.
    \item \textbf{Verification transcript:} allowed and denied actions per role.
    \item \textbf{Weekly audit query pack} against \texttt{ACCOUNT\_USAGE}.
\end{itemize}

\subsection*{Acceptance Criteria}
\begin{itemize}
    \item No privilege is granted directly to any user; no grant lands on \texttt{PUBLIC}.
    \item Access changes require editing only one layer (new table $\rightarrow$ future grant covers it; new hire $\rightarrow$ one functional grant).
    \item The audit query answers ``who can read this table'' including indirect inheritance.
    \item \texttt{ACCOUNTADMIN} and friends hold no business-object privileges.
\end{itemize}

\subsection*{Stretch Goals}
\begin{itemize}
    \item Add a network policy restricting service accounts to CI egress IPs and test from an off-list address.
    \item Write the offboarding script: revoke functional roles, transfer ownership, verify zero residual grants.
    \item Build the dormant-privilege report: grants with no matching \texttt{ACCESS\_HISTORY} usage in 90 days.
\end{itemize}

\section{Friday Real-World Architecture: 500 Users, 10 Business Units}
\index{Real-World Architecture!enterprise RBAC}
A multinational runs 500 Snowflake users across 10 business units (BUs), each with analysts, engineers, and pipeline accounts, plus cross-BU shared data products. Compliance demands least privilege, separation of duties, and quarterly access recertification.

\subsection{The Hierarchy Pattern}
Access roles are per-domain-and-layer: \texttt{AR\_SALES\_GOLD\_R}, \texttt{AR\_FIN\_SILVER\_W}, and so on. Functional roles are per-BU-and-job: \texttt{FR\_SALES\_ANALYST}, \texttt{FR\_FIN\_ENGINEER}. Shared data products get their own access roles granted across BUs explicitly---cross-BU visibility is a decision, not a default. Administration splits across \texttt{USERADMIN} (identity team), \texttt{SECURITYADMIN} (security team), and \texttt{SYSADMIN} (platform team); \texttt{ACCOUNTADMIN} is two named individuals, MFA-enforced, usage-alerted.

\begin{figure}[H]
    \centering
    \resizebox{\textwidth}{!}{%
    \begin{tikzpicture}[
        sys/.style={rectangle, rounded corners, draw=red!60!black, thick, fill=red!6, align=center, minimum width=2.6cm, minimum height=0.9cm, font=\scriptsize\bfseries},
        fr/.style={rectangle, rounded corners, draw=primary, thick, fill=primary!6, align=center, minimum width=2.6cm, minimum height=0.9cm, font=\scriptsize\bfseries},
        ar/.style={rectangle, rounded corners, draw=codegreen!60!black, thick, fill=green!8, align=center, minimum width=2.6cm, minimum height=0.9cm, font=\scriptsize\bfseries},
        usr/.style={rectangle, rounded corners, draw=codegray, thick, fill=white, align=center, minimum width=2.6cm, minimum height=0.9cm, font=\scriptsize},
        arrow/.style={-{Stealth[length=2.2mm]}, draw=primary, thick}
    ]
        \node[sys] (acct) at (6.5, 0) {ACCOUNTADMIN\\(break-glass, MFA)};
        \node[sys] (sec) at (2.8, -1.6) {SECURITYADMIN /\\USERADMIN};
        \node[sys] (sys) at (10.2, -1.6) {SYSADMIN};
        \node[fr] (fr1) at (1.2, -3.4) {FR\_SALES\_ANALYST};
        \node[fr] (fr2) at (4.4, -3.4) {FR\_FIN\_ENGINEER};
        \node[fr] (fr3) at (8.6, -3.4) {FR\_PIPELINE\_RUNNER};
        \node[ar] (ar1) at (1.2, -5.2) {AR\_SALES\_GOLD\_R};
        \node[ar] (ar2) at (4.4, -5.2) {AR\_FIN\_SILVER\_W};
        \node[ar] (ar3) at (8.6, -5.2) {AR\_SHARED\_PRODUCTS\_R};
        \node[usr] (u1) at (1.2, -7.0) {500 users\\(SSO + MFA)};
        \draw[arrow] (sec) -- (fr1);
        \draw[arrow] (sec) -- (fr2);
        \draw[arrow] (sec) -- (fr3);
        \draw[arrow] (fr1) -- (ar1);
        \draw[arrow] (fr2) -- (ar2);
        \draw[arrow] (fr2) -- (ar3);
        \draw[arrow] (fr3) -- (ar3);
        \draw[arrow] (u1) -- (fr1);
    \end{tikzpicture}%
    }
    \caption{Enterprise RBAC: users hold functional roles; functional roles roll up per-domain access roles; administration stays separated.}
    \label{fig:chapter17-enterprise-rbac}
\end{figure}

\subsection{Operating the Hierarchy}
Recertification is a report, not a meeting: \texttt{GRANTS\_TO\_USERS} joined to \texttt{LOGIN\_HISTORY} flags dormant access; per-BU owners approve the diff quarterly. Joiner-mover-leaver is automated through the identity provider's SCIM integration so role membership follows HR truth. And the standing guardrail queries---grants to \texttt{PUBLIC}, direct user grants, \texttt{ACCOUNTADMIN} activity---run weekly with results to the security channel.

\begin{pitfall}
\textbf{Hierarchy drift is inevitable without negative tests.} Quarterly, run the denial suite: each functional role attempts the actions it must \emph{not} perform. A hierarchy that is never tested at the boundary is a hope, not a control.
\end{pitfall}

\section{Interview Q\&A}
\begin{enumerate}
    \item \textbf{Q: Which system role must never run daily business workloads?}\\
    A: \texttt{ACCOUNTADMIN}---it inherits everything, so using it daily maximizes blast radius and destroys attribution.
    \item \textbf{Q: Access roles versus functional roles?}\\
    A: Access roles own privileges on objects (READ\_GOLD); functional roles describe jobs (ANALYST) and aggregate access roles. Users hold functional roles only.
    \item \textbf{Q: Why grant privileges to roles and never directly to users?}\\
    A: So access review, offboarding, and recertification are graph operations on a small role set instead of per-user privilege archaeology.
    \item \textbf{Q: What do secondary roles change in a session?}\\
    A: They let the session exercise additional roles' privileges simultaneously; object creation still belongs to the primary role, which is why primary choice matters.
    \item \textbf{Q: Which ACCOUNT\_USAGE views audit who can access a table?}\\
    A: \texttt{GRANTS\_TO\_ROLES} and \texttt{GRANTS\_TO\_USERS} (expanded through role inheritance), with \texttt{ACCESS\_HISTORY} and \texttt{LOGIN\_HISTORY} for actual usage.
\end{enumerate}

\section{Further Reading}
Snowflake's access-control overview and network-policy documentation define the privilege and network models \cite{snowflakeAccessControlDocs,snowflakeNetworkPoliciesDocs}. Chapter 18 builds column- and row-level enforcement on this foundation, and Chapter 19 extends it across account boundaries.

\section*{Key Takeaways}
\begin{itemize}
    \item System roles encode separation of duties; \texttt{ACCOUNTADMIN} is break-glass only.
    \item Access roles own object privileges; functional roles describe jobs; users hold functional roles.
    \item Future grants make least privilege the low-toil path.
    \item Secondary roles broaden a session; primary role still owns created objects.
    \item Network policies, MFA/SSO for humans, and key-pair auth for services form the perimeter.
    \item Audit is a standing query pack over \texttt{ACCOUNT\_USAGE}, plus quarterly negative tests.
\end{itemize}

\section{Exercises}
\begin{enumerate}
    \item \textbf{(Conceptual)} Explain why \texttt{PUBLIC} must receive no privileges.
    \item \textbf{(Conceptual)} A table created under the wrong primary role breaks the hierarchy. Explain and repair.
    \item \textbf{(Hands-on)} Build the full two-layer hierarchy for a third BU and prove it with the denial suite.
    \item \textbf{(Hands-on)} Write the ``who can read table X'' query expanding three levels of role inheritance.
    \item \textbf{(Hands-on)} Implement the dormant-privilege report and recertification diff.
    \item \textbf{(Design)} Design joiner-mover-leaver automation driven by the IdP, including the mover's access-diff.
    \item \textbf{(Design)} Specify network policies for humans, pipelines, and break-glass accounts.
    \item \textbf{(Operations)} Write the weekly guardrail queries: PUBLIC grants, direct user grants, ACCOUNTADMIN activity.
    \item \textbf{(Security)} A pipeline key leaked to a public repo. Write the incident runbook: revoke, rotate, audit blast radius.
    \item \textbf{(Interview)} Explain to an auditor why role-based grants are more reviewable than user-based grants.
\end{enumerate}
